\section{Spherical digit bounds over unrestricted integers}
\label{sec:signed}
\subsection{The exact inherited source and sign boundary}
All coordinates in this section range over $\Z$. The source equations are
the authenticated Round~17 centered dimension, grouped and sparse endpoints.
Their literal residual arrays and integral coordinate translations are
retained in \file{source_inputs}, bound by \file{SOURCE_PINS.json}.
The new changes are spherical batching of the eleven digit bounds, an
optional two-wire merger of adjacent denominator levels, and an optional
odd-square radix replacing one initial sign guard. All other rows remain
unchanged.

The inherited proof chain is \file{STRICT_BUDGET_PROOF.md},
\file{ORIENTATION_FIRST_PROOF.md}, \file{PROOF_315.md}, and the retained
Round~16 arithmetic-front proof, followed by the named C222/ZF/ZFC
calibration. In particular, the preceding signed frame and semantic
anchor establish the canonical orientation of both coordinates of every
principal Pell pair before using any of the eleven digit guards. This
orientation-first fact is indispensable; the present construction is not
an unrestricted reinterpretation of a natural system.

Let $t,F,V,B,x,y,N,q$ denote the semantic values before the integral
centering translations. The unchanged front has
\begin{align}
 t&=16+h_0^2,& V&=(F^2+1)^t,\\
 FB&=V+s_0,& B&=r_0+s_0+1,\\
 y&=V(V-1),& x&=F+1+(y+1)u_0,\\
 x&=1+[1+(t+1)y]u_t,& H_0&=1+\sum_{i=1}^{9}a_i,\\
 N&=B^x,& q&=(BH_0)^x.
\end{align}
Here the exponentiation equations indicate the interpretation of unchanged
polynomial certificates, not primitives of the final polynomial.
The factorial remainder and slack $r_0,s_0$ have retained nonnegativity
guards; the other symbols initially have only their declared signed
coordinate domains.

The square frame is $W=2+\sum(\text{power atom})^2\ge258$ and yields
$|A|>W^W$. All requested signed exponents have absolute value below $W$.
The squared semantic anchor bounds every classified Pell index. The
orientation pool, its grouped form, or its separated sparse residues then
recover $X_k(A),Y_k(A)$, allowing signed $k$ and signed $A$.
None of these arguments uses the eleven interval guards or the radix
output-slack guard changed here.

Consequently the inherited signed power lemma is available: for a base
$n\ge2$, a framed signed index $k$ and proposed integer output $m$, the
exact power certificate forces $k\ge0$ and $m=n^k$. If $k<0$, multiplying
the Pell congruence by $n^{-k}$ and using the large modulus bound would
force $mn^{-k}=1$, impossible over the integers. For base one and known
$k\ge0$, the request gives $m=1$. No output sign is presupposed.

In the baseline spherical variants, the old initial guard $q-B-1\ge0$
is retained. The inherited pre-factorial schedule therefore gives
\begin{equation}\label{eq:signedpreconditions}
 F\ge1,\quad B\ge2,\quad H_0\ge1,\quad y>F+2,\quad
 x\ge2+(t+1)y,\quad N=B^x,\quad q=(BH_0)^x.
\end{equation}
In particular $x\ge6$ and $N\ge16$. The odd-square variant obtains the
same conclusions by a separate acyclic schedule below.

\subsection{The old digit equations and their spherical replacement}
Write $D_k=q^k$. The exact factorial-binomial remainder satisfies
\begin{equation}\label{eq:signedfactorialdigit}
 (q+1)^{F^2+t}=(B+qH)D_t+r_f,
\end{equation}
where $H\in\Z$ is a high quotient. Each of the ten range remainders has
\begin{equation}\label{eq:signedrangedigit}
 (q+1)^a=ND_kQ_a+r_a,\qquad Q_a\in\Z.
\end{equation}
The groups are
\begin{center}\small
\begin{tabular}{@{}lll@{}}\toprule
Denominator & Remainders & Number\\\midrule
$D_t$ & factorial remainder and encoded index & 2\\
$D_{t+1}$ & encoded $a,b,c,d$ and derived label $h$ & 5\\
$D_{y+1}$ & the two stack representatives & 2\\
$D_x$ & the two strict quotient representatives & 2\\\bottomrule
\end{tabular}
\end{center}
The old quadratic interval guards are
\begin{equation}\label{eq:oldinterval}
 r(D-1-r)=a^2+b^2+c(c+1).
\end{equation}
For $D>0$, their right-hand sides are nonnegative, so they force
$0\le r<D$. Their exact source labels and centered coordinate identities
are checked in the finite transformation certificate.

For a group of remainders with common denominator $D$, replace all its
old guards by the single quadratic row
\begin{equation}\label{eq:sphere}
 \boxed{D^2-2\sum_i r_i^2=u^2+v^2+w(w+1).}
\end{equation}
All three new coordinates are unrestricted. The right-hand side is
nonnegative at every integer tuple. It also represents every nonnegative
integer: by the three-square theorem, $4n+1$ is a sum of three squares;
modulo four exactly one is odd. Writing the odd square as $(2w+1)^2$
and the other two as $(2u)^2,(2v)^2$ gives
$n=u^2+v^2+w(w+1)$.

Equation~\eqref{eq:sphere} gives, for each member independently,
\begin{equation}\label{eq:sphereradius}
 |r_i|\le D/\sqrt2<3D/4.
\end{equation}
It does not assume or immediately imply that $r_i$ is nonnegative. The
next argument supplies that conclusion using the exact digit equations.

\subsection{Signed uniqueness of a small low tail}
\begin{lemma}\label{lem:signedtail}
Let $q\ge2$, $k\ge1$, $D=q^k$, $a\ge0$, and
\[
 T(a,k)=\sum_{j=0}^{k-1}\binom ajq^j.
\]
If $T(a,k)<D/8$, $r\equiv T(a,k)\pmod D$, and $|r|<3D/4$, then
$r=T(a,k)$.
\end{lemma}
\begin{proof}
The integer $r-T(a,k)$ is a multiple of $D$ and satisfies
\[
 |r-T(a,k)|\le |r|+T(a,k)<7D/8<D.
\]
It must be zero.
\end{proof}
Binomial expansion gives $(q+1)^a\equiv T(a,k)\pmod D$.
Thus the literal rows \eqref{eq:signedfactorialdigit} and
\eqref{eq:signedrangedigit} provide the required congruence independently
of any sign on their quotients. The lemma explicitly excludes negative
noncanonical representatives.

\subsection{Recovering the factorial digit without circularity}
For $M\ge1$, $t\ge0$ and $0\le j\le t$,
\begin{equation}\label{eq:pascalbound}
 \binom{M+t}{j}\le(M+1)^t.
\end{equation}
For $j=t$, the product formula bounds each factor $(M+i)/i$ by $M+1$.
For $j<t$, induction on $t$ and Pascal's identity bound the two summands
by $2(M+1)^{t-1}\le(M+1)^t$; the boundary cases are immediate.

The unchanged $t$ request at base $F^2+1\ge1$ first yields
$V=(F^2+1)^t\ge1$. The retained ceiling relations and $B\ge1$ force
$F\ge1$. Since $FB\ge V\ge F^2+1$, we have $B>F$, hence
\[
 V\le FB<B^2.
\]
The endpoint gives $x\ge6$, and the exact radix gives $q\ge B^x$.
Therefore
\begin{equation}\label{eq:factorialsmall}
 B^2\le q/16,\qquad V<q/16.
\end{equation}
Applying \eqref{eq:pascalbound} with $M=F^2$, every coefficient through
position $t$ is at most $V$. Consequently
\begin{equation}\label{eq:factorialtail}
 T(F^2+t,t)<\frac{VD_t}{q-1}
 <\frac{qD_t}{16(q-1)}\le D_t/8.
\end{equation}
No factorial identity has been assumed. The $D_t$ sphere bounds $r_f$
even while its other coordinate remains uninterpreted, so
Lemma~\ref{lem:signedtail} gives $r_f=T(F^2+t,t)$.
Subtract this tail from \eqref{eq:signedfactorialdigit} and reduce modulo
$qD_t$ to obtain
\[
 B\equiv\binom{F^2+t}{t}\pmod q.
\]
The old output guard gives $0\le B<q$ in the baseline construction;
the odd-square schedule below proves $B<q$ directly. The binomial lies in
$[0,q)$ by \eqref{eq:factorialsmall}. Thus
$B=\binom{F^2+t}{t}$. Only now is the inherited square-radius factorial
fixed-point theorem applied to the unchanged ceiling inequalities,
yielding $F=t!$ and the exact reciprocal-CRT product $B$.

\subsection{Recovering the other ten intervals in order}
First handle ordinary encoded representatives, including the encoded-index
representative $C_0$. They are already nonnegative, and $H_0\ge1$.
For every ordinary representative $a$ and requested position
$0\le j\le k\le x$,
\[
 \binom aj\le\max(1,a)^k\le H_0^x=q/N.
\]
All four lower indices are at most $x$, by
$x\ge2+(t+1)y$ and $y\ge1$. Since $N\ge16$,
\begin{equation}\label{eq:ordinarytail}
 T(a,k)<\frac{qD_k}{N(q-1)}\le D_k/8.
\end{equation}
The spheres and Lemma~\ref{lem:signedtail} recover these ordinary tails,
in particular the tail for $C_0$. Independently, the retained literal
reciprocal-index congruence row gives $B\mid C_0+F^2+1$.
Since the now exact factorial product satisfies
$B>F^2+1$, $C_0=0$ is impossible. Thus $C_0\ge1$ and $H_0\ge2$.
The raw congruence does not itself depend on a remainder interval; stating
the ordinary tails first also makes clear that none uses the stronger
derived-label majorant.

For the derived label, the unchanged constructor ledger gives
\[
 h\le H_0^\ell,\qquad \ell=4194304.
\]
The inherited factorial estimate $\ell(t+1)<2F$ and endpoint bound
$2F<x$ imply $\ell(t+1)<x$. Its needed binomial coefficients through
position $t+1$ are therefore also at most $H_0^x=q/N$.
The same strict tail estimate identifies its remainder. All eleven old
interval conclusions have now been recovered.

More explicitly, after subtracting a range tail, divide
\eqref{eq:signedrangedigit} by $D_k$. The expansion becomes
\[
 NQ_a=\binom ak+qH_a
\]
for an integer $H_a$. Since $N\mid q$, this yields
$N\mid\binom ak$, exactly the inherited range predicate.
The unchanged prime-power precision, residue extraction and local Horn
flow arguments then yield the original finite-refutation interpretation.
No signed quotient has been treated as natural before this recovery.

\subsection{Both witness directions and adjacent-level merging}
For completeness, start with a finite refutation and construct the inherited
cofinal semantic witnesses. Choose all digit remainders canonically as
$T(a,k)$. By \eqref{eq:factorialtail} and \eqref{eq:ordinarytail}, each is
less than its denominator divided by eight. In a group with $m\le5$,
\[
 2\sum_iT_i^2<\frac{m}{32}D^2<D^2.
\]
The integer $D^2-2\sum_iT_i^2$ is nonnegative, so its ternary representation
supplies the new guard coordinates. All unchanged residuals already vanish.
Conversely, the preceding soundness proof restores every old interval from
a new zero. Each $r(D-1-r)\ge0$ then has a ternary representation, filling
the deleted old coordinates afresh. Undoing the recorded integral
translation gives an actual zero of the calibrated source. This proves
integer existential equivalence, without claiming that arbitrary old
private witnesses extend unchanged.

For the optional merger, introduce two explicit product coordinates
\[
 z_f=qr_f,\qquad z_C=qr_C.
\]
Replace the first two spheres by one at denominator $D_{t+1}=qD_t$,
using $z_f,z_C$ and the five old $D_{t+1}$ remainders. The new sphere gives
$|qr_f|,|qr_C|<3qD_t/4$, hence the old radius bound because $q>0$.
The soundness argument is unchanged. Canonical scaled tails are below
$D_{t+1}/8$, and now
\[
 2\sum_{i=1}^{7}T_i^2<\frac7{32}D_{t+1}^2<D_{t+1}^2,
\]
so completeness still holds. Two ternary triples become one, saving three
coordinates, while two product wires are added: net one fewer coordinate
and one additional residual. Every residual remains quadratic.

\subsection{An odd square replaces the initial radix sign guard}
The separately identified odd-radix variants remove the three-coordinate
ternary guard for $q-B-1\ge0$ and add one unrestricted $z$ with
\begin{equation}\label{eq:oddradix}
 q=(2z+1)^2.
\end{equation}
This supplies only $q\ge1$ initially. The following schedule replaces the
old use of $q\ge2$.
\begin{enumerate}
\item The unchanged frame, anchor and orientation proof gives canonical
principal Pell pairs without the removed guard. Since $t\ge16$, the
request at base $F^2+1$ yields $V=(F^2+1)^t$, including the base-one case.
The retained ceiling guards then give $F\ge1$, $B\ge2$, $V<B^2$ and
$y>F+2$.
\item Now $q+1\ge2$. Each actual encoded representative is the exponent
of a numerator request at this base, so the signed power lemma makes every
such representative nonnegative. Hence $H_0\ge1$. The derived-label
request similarly gives a nonnegative index, with zero already excluded
by the squared anchor.
\item Since $BH_0\ge2$, its existing $x$ request forces $x\ge0$ and
$q=(BH_0)^x$.
\item In $x=F+1+(y+1)u_0$, a value $u_0\le-1$ would imply
$x\le F-y<0$. Thus $x\ge F+1\ge2$. The other endpoint forces
$u_t\ge1$, so $x\ge2+(t+1)y$. Consequently
$q\ge B^x\ge B^2>B$ and $q\ge4$.
\item The remaining powers, including $N=B^x$ and every denominator,
are now exact. The missing output slack $q-B-1$ is nonnegative before
factorial decoding. The preceding spherical proof applies in its stated
order.
\end{enumerate}
No removed guard is used to establish its own replacement.

For completeness, choose the beta code $x$ even in the inherited cofinal
CRT construction. Since $y=V(V-1)$ is even, every beta modulus
$1+(i+1)y$ is odd, and adding their odd common CRT modulus changes parity.
A sufficiently large even representative preserves all encoded values and
nonnegative endpoint quotients. After $F=t!$, each factor of $B$ is odd,
so $N^2$ is odd. Every later encoded representative can be changed by
$N^2$ while preserving its required congruences. Choose $a,b$ of opposite
parity, as already needed for even semantic-anchor index $h$. Independently
choose one of the other seven parities so the sum of all nine
representatives is even; then $H_0$ is odd.

The dominant $a$ can still lie in $[2N^2,4N^2)$ with its prescribed parity,
and the other representatives remain at most $2N^2$. Thus the inherited
anchor dominance and cap estimates remain valid. Now $BH_0$ is odd and
$x$ even, so the positive integer $R=(BH_0)^{x/2}$ is odd. Taking
$z=(R-1)/2$ gives \eqref{eq:oddradix}. All other CRT, radix and Pell
witness choices proceed unchanged. This proves that the square restriction
is compatible with completeness, not merely sufficient for soundness.

\subsection{What the finite checker establishes}
The producer reads only copied, hash-bound source equations and their
coordinate translations. The separately implemented checker imports no
producer. It reconstructs each replacement in named coordinates, verifies
every untouched residual, expands each square and collects all coefficients,
and checks liveness, degree, support and height. Deliberate changes to sphere,
frame and odd-square rows must be rejected. Finite Pascal or tail examples
serve as diagnostics only; the preceding unbounded arguments prove the
mathematical implication. Together with the retained ordinary calibration,
this establishes \eqref{eq:calibration} for each selected unrestricted
quartic.
